% Image credits for the photographs used in this volume (Book 2 only).
% Machine-readable license records live in images/book2/CREDITS.md (and
% images/book1/CREDITS.md for the three photographs reused from Book 1);
% keep the files in sync when adding or removing a photograph.
\chapter*{Créditos das imagens}

% Under bidi=basic a Latin run is ONE directional box, so a long author or
% institution name cannot be broken at all; \allowbreak inside it does nothing.
% The Arabic edition reported this page as the book's only overfull boxes.
% \sloppy must be in force when the paragraph is BROKEN, hence the \par before
% \endgroup -- without it the group closes first and the tolerance is restored.
\begingroup\sloppy

Os desenhos deste livro são originais. As fotografias são usadas sob
as respectivas licenças livres, com agradecimento aos seus autores:

\begin{itemize}
  \item \emph{Gregor Mendel} (capítulo sobre a meiose): fotografia,
        século \textsc{xix}, Wikimedia Commons, domínio público.
  \item \emph{Alexander Fleming} (capítulo sobre a resistência aos
        antibióticos): fotografia, Imperial War Museums, 1943, via
        Wikimedia Commons, domínio público.
  \item \emph{Brotamento do HIV em um linfócito} (capítulo sobre a
        imunidade adaptativa): micrografia eletrônica de
        C.~Goldsmith, Centers for Disease Control and Prevention, via
        Wikimedia Commons, domínio público.
  \item \emph{Os tentilhões de Darwin} (capítulo sobre a
        diversificação): gravura de John Gould, 1845, Wikimedia
        Commons, domínio público.
  \item \emph{Homúnculo sensorial} (capítulo sobre o movimento):
        ilustração de OpenStax \emph{Anatomy and Physiology}, via
        Wikimedia Commons, CC~BY~3.0.
  \item \emph{Charles Darwin} (capítulo sobre as árvores
        filogenéticas): fotografia de Henry Maull e John Fox, por
        volta de 1854, Wikimedia Commons, domínio público.
  \item \emph{Edward Jenner} (capítulo sobre a imunidade adaptativa):
        pintura a óleo, Wellcome Collection via Wikimedia Commons,
        CC~BY~4.0.
  \item \emph{Fóssil de Archaeopteryx} (capítulo sobre a ancestralidade
        comum): fotografia de James L.~Amos, Wikimedia Commons, CC0.
\end{itemize}

Os links completos das fontes e os endereços das licenças de cada
fotografia estão em \texttt{images/\allowbreak book2/\allowbreak CREDITS.md} e
\texttt{images/\allowbreak book1/\allowbreak CREDITS.md}, no repositório do livro.

\bigskip
Além das fotografias e dos desenhos originais, cerca de cinquenta
figuras espalhadas pelos capítulos do volume são ilustrações geradas
por inteligência artificial, produzidas com a geração de imagens da
OpenAI a partir de instruções escritas pelos editores do livro, e
revisadas quanto à exatidão biológica antes da inclusão. A instrução
completa de cada imagem gerada está registrada em
\texttt{images/\allowbreak book2/\allowbreak ai/\allowbreak PROMPTS.md}, no repositório.
\par\endgroup
\clearpage
